\documentclass[a4paper,12pt]{article}

% ========================================================================
% Pacotes
% ========================================================================
\usepackage[margin=2.5cm]{geometry}
\usepackage[utf8]{inputenc}
\usepackage[T1]{fontenc}
\usepackage[brazil]{babel}
\usepackage{tikz}
\usepackage{amsmath}
\usepackage{amssymb}
\usepackage{graphicx}
\usepackage{float}
\usepackage{listings}
\usepackage{xcolor}
\usepackage{caption}
\usepackage{subcaption}
\usepackage{fancyhdr}
\usepackage{titlesec}

% TikZ
\usetikzlibrary{
    arrows.meta, positioning, calc, shapes.geometric,
    fit, decorations.markings, patterns
}

% ========================================================================
% Estilos TikZ compartilhados
% ========================================================================
\tikzset{
    block/.style = {
        rectangle, draw, very thick,
        minimum width=3.5cm, minimum height=1.8cm,
        font=\sffamily\small, align=center, text width=3cm
    },
    wide block/.style = {
        block, minimum width=5cm, text width=4.5cm
    },
    small block/.style = {
        rectangle, draw, thick,
        minimum width=2cm, minimum height=1cm,
        font=\sffamily\footnotesize, align=center
    },
    operator/.style = {
        circle, draw, very thick,
        minimum size=9mm, font=\sffamily\large,
        inner sep=0pt, fill=white
    },
    dot/.style = {
        circle, fill=black, minimum size=3pt, inner sep=0pt
    },
    port/.style = {font=\sffamily\bfseries\small},
    signal/.style = {font=\sffamily\footnotesize, fill=white, inner sep=1pt},
    bus/.style = {font=\sffamily\scriptsize},
    entity border/.style = {
        rectangle, draw, dashed, thick,
        rounded corners=3pt, inner sep=8pt
    },
    register/.style = {
        rectangle, draw, thick,
        minimum width=1.0cm, minimum height=0.8cm,
        font=\sffamily\footnotesize, align=center
    },
    data arrow/.style = {-Stealth, thick},
    ctrl arrow/.style = {-Stealth, thin, dashed},
    const/.style = {font=\sffamily\itshape\footnotesize},
    annotation/.style = {font=\sffamily\small, align=left},
    stage sep/.style = {dashed, thin, black!50},
    >=Stealth, thick, font=\sffamily
}

% ========================================================================
% Estilo VHDL para listings
% ========================================================================
\definecolor{vhdlkeyword}{RGB}{0, 0, 160}
\definecolor{vhdlcomment}{RGB}{0, 128, 0}
\definecolor{vhdlstring}{RGB}{160, 0, 0}
\definecolor{vhdlbg}{RGB}{248, 248, 248}
\definecolor{vhdlnumber}{RGB}{100, 100, 100}

\lstdefinestyle{vhdlstyle}{
    language=VHDL,
    backgroundcolor=\color{vhdlbg},
    basicstyle=\ttfamily\scriptsize,
    keywordstyle=\color{vhdlkeyword}\bfseries,
    commentstyle=\color{vhdlcomment}\itshape,
    stringstyle=\color{vhdlstring},
    numberstyle=\tiny\color{vhdlnumber},
    numbers=left,
    numbersep=5pt,
    frame=single,
    rulecolor=\color{black!30},
    breaklines=true,
    breakatwhitespace=false,
    tabsize=4,
    showstringspaces=false,
    captionpos=b,
    xleftmargin=2em,
    framexleftmargin=1.5em,
    morekeywords={data_t, data_long_t, data_array_t, signed, std_logic, std_logic_vector}
}

\lstdefinestyle{logstyle}{
    backgroundcolor=\color{black!5},
    basicstyle=\ttfamily\scriptsize,
    frame=single,
    rulecolor=\color{black!30},
    breaklines=true,
    captionpos=b,
    xleftmargin=1em,
    framexleftmargin=0.5em,
}

% ========================================================================
% Cabecalho
% ========================================================================
\pagestyle{fancy}
\fancyhf{}
\fancyhead[L]{\sffamily\small Rede Neural DPD em VHDL}
\fancyhead[R]{\sffamily\small Etapa 1: Input Processing Block}
\fancyfoot[C]{\thepage}

% ========================================================================
% Documento
% ========================================================================
\begin{document}

% -----------------------------------------------------------------------
% Capa
% -----------------------------------------------------------------------
\begin{titlepage}
\centering
\vspace*{3cm}

{\huge\sffamily\bfseries Implementa\c{c}\~{a}o em VHDL de\\[5pt] Rede Neural para DPD}\\[1cm]
{\Large\sffamily Modelo CMEMP -- Ponto Fixo Q16.16}\\[2cm]

{\large\sffamily\bfseries Relat\'{o}rio de Progresso -- Etapa 1}\\[0.5cm]
{\large\sffamily Bloco de Processamento de Entrada}\\[3cm]

\vfill

{\sffamily\today}

\end{titlepage}

% -----------------------------------------------------------------------
% Sum\'{a}rio
% -----------------------------------------------------------------------
\tableofcontents
\newpage

% =======================================================================
\section{Vis\~{a}o Geral do Sistema}
% =======================================================================

O objetivo deste projeto \'{e} implementar em VHDL uma rede neural para lineariza\c{c}\~{a}o de amplificadores de pot\^{e}ncia (DPD -- \textit{Digital Pre-Distortion}), utilizando o modelo CMEMP (\textit{Cross-Memory Even-order Modified Parallel}). A rede opera em aritm\'{e}tica de ponto fixo Q16.16 (32 bits com sinal, 16 bits fracion\'{a}rios).

A arquitetura completa \'{e} dividida em tr\^{e}s est\'{a}gios de processamento sequencial, conforme ilustrado na Figura~\ref{fig:sistema_completo}:

\begin{enumerate}
    \item \textbf{Est\'{a}gio 1 -- Processamento de Entrada:} Converte as amostras IQ retangulares para coordenadas polares, realiza decomposi\c{c}\~{a}o trigonom\'{e}trica e mant\'{e}m buffers deslizantes com hist\'{o}rico de $M+1$ amostras.
    \item \textbf{Est\'{a}gio 2 -- Camada Convolucional:} Aplica 16 filtros convolucionais 1D (serializados) sobre os tr\^{e}s canais (envelope, cosseno e seno), seguidos de ativa\c{c}\~{a}o ReLU e \textit{average pooling}.
    \item \textbf{Est\'{a}gio 3 -- MLP e Sa\'{i}da:} Duas redes MLP paralelas (ramo real e imagin\'{a}rio) processam o vetor de caracter\'{i}sticas, seguidas de rota\c{c}\~{a}o de fase via multiplica\c{c}\~{a}o complexa.
\end{enumerate}

O plano de implementa\c{c}\~{a}o segue uma abordagem \textit{bottom-up}: construir, testar e validar cada bloco individualmente antes de integr\'{a}-los. \textbf{Este relat\'{o}rio documenta a primeira etapa: o projeto e valida\c{c}\~{a}o do \texttt{Input\_Processing\_Block}}, que \'{e} o bloco fundamental do Est\'{a}gio~1.

% -----------------------------------------------------------------------
% Diagrama do Sistema Completo (em portugu\^{e}s, sem FSM)
% -----------------------------------------------------------------------
\begin{figure}[H]
\centering
\resizebox{\textwidth}{!}{%
\begin{tikzpicture}[
    node distance=1.5cm and 2.5cm,
    >=Stealth, thick, font=\sffamily
]

\tikzset{
    io_node/.style={font=\sffamily\bfseries},
    dot/.style={circle, fill=black, inner sep=0pt, minimum size=4pt}
}

% --- Tr\^{e}s est\'{a}gios ---
\node[wide block, minimum height=4.5cm, text width=4.0cm] (S1) at (0, 0) {%
    \textbf{Est\'{a}gio 1}\\
    \textbf{Processamento}\\
    \textbf{de Entrada}\\[5pt]
    {\footnotesize Top\_NeuralNet\_S1}\\[3pt]
    {\footnotesize IQ $\rightarrow$ Polar}\\
    {\footnotesize Decomp. trigonom.}\\
    {\footnotesize Buffers deslizantes}
};

\node[wide block, minimum height=4.5cm, text width=4.0cm] (S2) at ($(S1.east) + (6.5, 0)$) {%
    \textbf{Est\'{a}gio 2}\\
    \textbf{Convolu\c{c}\~{a}o}\\[5pt]
    {\footnotesize Conv\_Layer\_Serial}\\[3pt]
    {\footnotesize Conv1D 3 canais}\\
    {\footnotesize 16 filtros (serial.)}\\
    {\footnotesize ReLU + Avg Pool}
};

\node[wide block, minimum height=4.5cm, text width=4.0cm] (S3) at ($(S2.east) + (6.5, 0)$) {%
    \textbf{Est\'{a}gio 3}\\
    \textbf{MLP + Sa\'{i}da}\\[5pt]
    {\footnotesize BackEnd\_Wrapper}\\[3pt]
    {\footnotesize 2 ramos MLP}\\
    {\footnotesize Rota\c{c}\~{a}o de fase}\\
    {\footnotesize Mult. complexa}
};

% --- Entradas ---
\node[io_node] (Iin) at ($(S1.west) + (-3.5, 1.0)$) {$i_{in}$ \normalfont\footnotesize [data\_t]};
\node[io_node] (Qin) at ($(S1.west) + (-3.5, -0.3)$) {$q_{in}$ \normalfont\footnotesize [data\_t]};
\node[io_node] (Vin) at ($(S1.west) + (-3.5, -1.5)$) {input\_valid};

\draw[data arrow] (Iin) -- ($(S1.west) + (0, 1.0)$);
\draw[data arrow] (Qin) -- ($(S1.west) + (0, -0.3)$);
\draw[ctrl arrow] (Vin) -- ($(S1.west) + (0, -1.5)$);

% --- S1 -> S2: barramentos de buffer ---
\draw[data arrow] ($(S1.east) + (0, 1.2)$) -- ($(S2.west) + (0, 1.2)$)
    node[signal, midway, above] {buf\_env$[0\!:\!5]$}
    node[bus, midway, below] {$/6$};

\draw[data arrow] ($(S1.east) + (0, 0.0)$) -- ($(S2.west) + (0, 0.0)$)
    node[signal, midway, above] {buf\_cos$[0\!:\!5]$}
    node[bus, midway, below] {$/6$};

\draw[data arrow] ($(S1.east) + (0, -1.2)$) -- ($(S2.west) + (0, -1.2)$)
    node[signal, midway, above] {buf\_sin$[0\!:\!5]$}
    node[bus, midway, below] {$/6$};

% --- S1 -> S3: bypass theta ---
\coordinate (thetaStart) at ($(S1.east) + (0, -1.8)$);
\coordinate (thetaEnd) at ($(S3.west) + (0, -1.8)$);
\coordinate (thetaBelow) at ($(S2.south) + (0, -1.2)$);

\draw[data arrow] (thetaStart) -- ++(0, -1.5) -- (thetaBelow);
\node[register, fill=white] (ThetaReg) at (thetaBelow) {$\theta_{cap}$};
\node[signal, below=0.2cm of ThetaReg] {\scriptsize theta\_captured};
\draw[data arrow] (ThetaReg.east) -| (thetaEnd);
\node[signal] at ($(thetaStart) + (0.3, -0.8)$) {\scriptsize $\theta$};

% --- S2 -> S3: vetor de features ---
\draw[data arrow] ($(S2.east) + (0, 0.0)$) -- ($(S3.west) + (0, 0.0)$)
    node[signal, midway, above] {x\_flat$[0\!:\!15]$}
    node[bus, midway, below] {$/16$};

% --- Sa\'{i}das ---
\node[io_node] (Iout) at ($(S3.east) + (3.0, 0.8)$) {$i_{out}$ \normalfont\footnotesize [data\_t]};
\node[io_node] (Qout) at ($(S3.east) + (3.0, -0.5)$) {$q_{out}$ \normalfont\footnotesize [data\_t]};
\node[io_node] (Vout) at ($(S3.east) + (3.0, -1.5)$) {output\_valid};

\draw[data arrow] ($(S3.east) + (0, 0.8)$) -- (Iout);
\draw[data arrow] ($(S3.east) + (0, -0.5)$) -- (Qout);
\draw[ctrl arrow] ($(S3.east) + (0, -1.5)$) -- (Vout);

% --- Anota\c{c}\~{a}o ---
\node[annotation] at ($(S2.south) + (0, -3.0)$) {%
    {\footnotesize Formato: Q16.16 (32 bits com sinal, 16 bits fracion\'{a}rios) \quad $M\!=\!5$, $N_{\text{filtros}}\!=\!16$, $K\!=\!4$}
};

\end{tikzpicture}%
}
\caption{Arquitetura geral do sistema -- tr\^{e}s est\'{a}gios de processamento sequencial. A fase $\theta$ realiza \textit{bypass} do Est\'{a}gio~2 atrav\'{e}s de um registrador de captura.}
\label{fig:sistema_completo}
\end{figure}


% =======================================================================
\section{Est\'{a}gio 1: Processamento de Entrada}
% =======================================================================

O Est\'{a}gio~1 (\texttt{Top\_NeuralNet\_S1}) \'{e} respons\'{a}vel por converter o sinal IQ de entrada em uma representa\c{c}\~{a}o adequada para as camadas subsequentes da rede neural. A Figura~\ref{fig:estagio1} apresenta seu diagrama de blocos.

O fluxo de processamento \'{e}:
\begin{enumerate}
    \item O \texttt{Input\_Processing\_Block} recebe as componentes $i_{in}$ e $q_{in}$ e, utilizando o algoritmo CORDIC, converte-as em coordenadas polares: envelope $A(n)$ e fase $\theta(n)$.
    \item A fase normalizada $\Delta\theta$ alimenta duas LUTs trigonom\'{e}tricas que calculam $\cos(\Delta\theta)$ e $\sin(\Delta\theta)$ em 4 ciclos de pipeline.
    \item Um bloco de atraso de alinhamento (4 ciclos) compensa a lat\^{e}ncia das LUTs para o envelope e a fase.
    \item Tr\^{e}s registradores de deslocamento (\textit{Shift Registers}) mant\^{e}m o hist\'{o}rico de $M+1 = 6$ amostras para cada canal.
\end{enumerate}

% -----------------------------------------------------------------------
% Diagrama do Est\'{a}gio 1 (corrigido e em portugu\^{e}s)
% -----------------------------------------------------------------------
\begin{figure}[H]
\centering
\resizebox{\textwidth}{!}{%
\begin{tikzpicture}[
    node distance=1.2cm and 1.8cm,
    >=Stealth, thick, font=\sffamily
]

\tikzset{
    io_node/.style={font=\sffamily\bfseries},
    dot/.style={circle, fill=black, inner sep=0pt, minimum size=4pt}
}

% --- Bloco de Processamento de Entrada ---
\node[block, minimum width=4.0cm, minimum height=4.0cm, text width=3.5cm] (IPB) at (0, 0) {%
    \textbf{Input}\\
    \textbf{Processing}\\
    \textbf{Block}\\[3pt]
    {\footnotesize CORDIC + Phase}\\
    {\footnotesize Unwrap}\\
    {\footnotesize (40 ciclos)}
};

% --- Entradas ---
\node[io_node] (Iin) at (-4.5, 1.0) {$i_{in}$};
\node[io_node] (Qin) at (-4.5, -0.5) {$q_{in}$};
\node[io_node] (Vin) at (-4.5, -2.0) {input\_valid};

\draw[data arrow] (Iin) -- ($(IPB.west) + (0, 1.0)$);
\draw[data arrow] (Qin) -- ($(IPB.west) + (0, -0.2)$);
\draw[ctrl arrow] (Vin) -- ($(IPB.west) + (0, -1.4)$);

% --- Sa\'{i}das do IPB ---
\coordinate (envOut) at ($(IPB.east) + (0, 1.2)$);
\coordinate (thetaNormOut) at ($(IPB.east) + (0, 0.0)$);
\coordinate (phaseWOut) at ($(IPB.east) + (0, -1.2)$);
\coordinate (enRawOut) at ($(IPB.south) + (0.5, 0)$);

% --- LUTs trigonom\'{e}tricas ---
\node[block, minimum width=2.8cm, minimum height=1.5cm, text width=2.4cm] (LUT_C) at (5.5, 1.5) {%
    \textbf{LUT\_Trig}\\
    {\footnotesize (cosseno)}\\
    {\footnotesize $\cos(\Delta\theta)$}
};

\node[block, minimum width=2.8cm, minimum height=1.5cm, text width=2.4cm] (LUT_S) at (5.5, -1.5) {%
    \textbf{LUT\_Trig}\\
    {\footnotesize (seno)}\\
    {\footnotesize $\sin(\Delta\theta)$}
};

% theta_norm bifurca para ambas LUTs
\draw[data arrow] (thetaNormOut) -- ++(1.5, 0) coordinate (thetaFork);
\node[dot] at (thetaFork) {};
\draw[data arrow] (thetaFork) |- (LUT_C.west)
    node[signal, pos=0.3, above] {\scriptsize $\theta_{norm}$};
\draw[data arrow] (thetaFork) |- (LUT_S.west);

% --- Atraso de Alinhamento ---
\node[small block, minimum width=2.8cm, minimum height=1.8cm, text width=2.4cm] (Delay) at (5.5, 4.5) {%
    \textbf{Atraso de}\\
    \textbf{Alinhamento}\\
    {\scriptsize 4 ciclos (LUT)}
};

% envelope_s -> Atraso
\draw[data arrow] (envOut) -- ++(1.5, 0) |- ($(Delay.west) + (0, 0.3)$)
    node[signal, pos=0.2, above] {\scriptsize envelope\_s};

% phase_wrapped -> Atraso
\draw[data arrow] (phaseWOut) -- ++(0.8, 0) |- ($(Delay.west) + (0, -0.3)$)
    node[signal, pos=0.15, below] {\scriptsize phase\_wrapped};

% --- Atraso do Enable ---
\node[small block, minimum width=2.8cm, minimum height=0.8cm, text width=2.4cm] (EnDelay) at (5.5, -4.0) {%
    \textbf{Atraso Enable}\\
    {\scriptsize 4 ciclos}
};
\draw[ctrl arrow] (enRawOut) |- (EnDelay.west)
    node[signal, pos=0.3, below] {\scriptsize en\_raw};

% --- Registradores de Deslocamento ---
\node[block, minimum width=3.2cm, minimum height=1.3cm, text width=2.8cm] (SR_E) at (11.5, 4.5) {%
    \textbf{Shift\_Reg}\\
    {\footnotesize (envelope, $M\!=\!5$)}
};

\node[block, minimum width=3.2cm, minimum height=1.3cm, text width=2.8cm] (SR_C) at (11.5, 1.5) {%
    \textbf{Shift\_Reg}\\
    {\footnotesize (cosseno, $M\!=\!5$)}
};

\node[block, minimum width=3.2cm, minimum height=1.3cm, text width=2.8cm] (SR_S) at (11.5, -1.5) {%
    \textbf{Shift\_Reg}\\
    {\footnotesize (seno, $M\!=\!5$)}
};

% Atraso -> Shift_Reg (envelope)
\draw[data arrow] (Delay.east) -- ++(0.5, 0) |- (SR_E.west)
    node[signal, pos=0.4, above] {\scriptsize env\_aligned};

% LUT -> Shift_Reg
\draw[data arrow] (LUT_C.east) -- (SR_C.west)
    node[signal, midway, above] {\scriptsize cos\_dp};
\draw[data arrow] (LUT_S.east) -- (SR_S.west)
    node[signal, midway, above] {\scriptsize sin\_dp};

% Distribui\c{c}\~{a}o do enable
\coordinate (enAligned) at ($(EnDelay.east) + (1.5, 0)$);
\draw[ctrl arrow] (EnDelay.east) -- (enAligned);
\node[dot] at (enAligned) {};
\draw[ctrl arrow] (enAligned) |- ($(SR_E.south) + (-0.5, 0)$)
    node[signal, pos=0.1, right] {\scriptsize wr\_en};
\draw[ctrl arrow] (enAligned) |- ($(SR_C.south) + (-0.5, 0)$);
\draw[ctrl arrow] (enAligned) |- ($(SR_S.south) + (-0.5, 0)$);

% --- Sa\'{i}das (direita) ---
\node[io_node, right=1.5cm of SR_E] (OutE) {buf\_env$[0\!:\!5]$};
\node[io_node, right=1.5cm of SR_C] (OutC) {buf\_cos$[0\!:\!5]$};
\node[io_node, right=1.5cm of SR_S] (OutS) {buf\_sin$[0\!:\!5]$};

\draw[data arrow] (SR_E.east) -- (OutE) node[bus, midway, below] {$/6$};
\draw[data arrow] (SR_C.east) -- (OutC) node[bus, midway, below] {$/6$};
\draw[data arrow] (SR_S.east) -- (OutS) node[bus, midway, below] {$/6$};

% --- theta_unwrapped sa\'{i}da (CORRIGIDO: posicionado acima do SR_E) ---
\node[io_node] (OutTheta) at ($(SR_E.east) + (1.5, 1.5)$) {theta\_unwrapped\_out};
\draw[data arrow] ($(Delay.east) + (0.5, -0.3)$) -- ++(1.0, 0) |- (OutTheta)
    node[signal, pos=0.3, right] {\scriptsize $\theta_{aligned}$};

% --- stage_valid ---
\node[io_node] (OutValid) at ($(enAligned) + (4.5, 0)$) {stage\_valid};
\draw[ctrl arrow] (enAligned) -- ++(0.5, 0) |- (OutValid);

\end{tikzpicture}%
}
\caption{Diagrama de blocos do Est\'{a}gio 1 (\texttt{Top\_NeuralNet\_S1}): convers\~{a}o IQ$\,\rightarrow\,$polar, decomposi\c{c}\~{a}o trigonom\'{e}trica e buffers deslizantes. A lat\^{e}ncia das LUTs (4 ciclos) \'{e} compensada pelo bloco de atraso de alinhamento.}
\label{fig:estagio1}
\end{figure}


% =======================================================================
\section{Input Processing Block}
\label{sec:input_processing}
% =======================================================================

O \texttt{Input\_Processing\_Block} \'{e} o primeiro bloco a ser implementado e validado. Ele \'{e} respons\'{a}vel por:

\begin{itemize}
    \item \textbf{Convers\~{a}o retangular para polar} usando o algoritmo CORDIC (\textit{Coordinate Rotation Digital Computer}), com 40 ciclos de lat\^{e}ncia.
    \item \textbf{Corre\c{c}\~{a}o de ganho CORDIC}: a magnitude produzida pelo CORDIC possui ganho intr\'{i}nseco $K \approx 1{,}1644$. O bloco multiplica pelo inverso $K^{-1} \approx 0{,}8588$ e extrai os bits $[47\!:\!16]$ do produto de 64 bits.
    \item \textbf{Escala de fase}: a sa\'{i}da de fase do CORDIC est\'{a} em formato de alta resolu\c{c}\~{a}o e \'{e} convertida para Q16.16 atrav\'{e}s de \textit{shift right} de 13 bits.
    \item \textbf{L\'{o}gica de desembrulhamento de fase} (\textit{phase unwrapping}): detecta saltos de $2\pi$ e mant\'{e}m continuidade da fase acumulada $\theta_{unwrapped}(n)$.
\end{itemize}


\subsection{Interface do Bloco}

\begin{table}[H]
\centering
\caption{Portas do \texttt{Input\_Processing\_Block}.}
\label{tab:ipb_ports}
\small
\begin{tabular}{llll}
\hline
\textbf{Porta} & \textbf{Dir.} & \textbf{Tipo} & \textbf{Descri\c{c}\~{a}o} \\
\hline
\texttt{clk}                & in  & std\_logic & Rel\'{o}gio do sistema \\
\texttt{rst}                & in  & std\_logic & Reset s\'{i}ncrono \\
\texttt{i\_in}              & in  & data\_t    & Componente I (parte real) \\
\texttt{q\_in}              & in  & data\_t    & Componente Q (parte imagin\'{a}ria) \\
\texttt{input\_valid}       & in  & std\_logic & Indica amostra v\'{a}lida \\
\hline
\texttt{envelope\_out}      & out & data\_t    & Magnitude $A(n)$ corrigida \\
\texttt{theta\_unwrapped\_out} & out & data\_t & Fase acumulada (sem saltos) \\
\texttt{theta\_norm\_out}   & out & data\_t    & $\Delta\theta(n)$ normalizado $[-\pi, +\pi]$ \\
\texttt{phase\_wrapped\_out}& out & data\_t    & Fase direta $\theta_{raw}(n)$ \\
\texttt{buffers\_enable}    & out & std\_logic & Pulso de habilita\c{c}\~{a}o ap\'{o}s lat\^{e}ncia CORDIC \\
\hline
\end{tabular}
\end{table}


\subsection{C\'{o}digo VHDL}

O c\'{o}digo completo do \texttt{Input\_Processing\_Block} \'{e} apresentado a seguir:

\begin{lstlisting}[style=vhdlstyle, caption={C\'{o}digo VHDL do \texttt{Input\_Processing\_Block} (\texttt{input\_processing.vhd}).}, label={lst:ipb}]
library IEEE;
use IEEE.STD_LOGIC_1164.ALL;
use IEEE.NUMERIC_STD.ALL;
use work.pkg_neural_types.ALL;

entity Input_Processing_Block is
    Port (
        clk                : in  std_logic;
        rst                : in  std_logic;
        i_in               : in  data_t;
        q_in               : in  data_t;
        input_valid        : in  std_logic;
        envelope_out       : out data_t;
        theta_unwrapped_out: out data_t;
        theta_norm_out     : out data_t;
        phase_wrapped_out  : out data_t;
        buffers_enable     : out std_logic
    );
end Input_Processing_Block;

architecture Behavioral of Input_Processing_Block is

    component cordic_rect_to_polar
    port (
        clk : in std_logic;
        x_in : in std_logic_vector(31 downto 0);
        y_in : in std_logic_vector(31 downto 0);
        x_out : out std_logic_vector(31 downto 0);
        phase_out : out std_logic_vector(31 downto 0)
    );
    end component;

    constant PI_FIXED : data_t := to_signed(205887, 32);
    constant TWO_PI   : data_t := to_signed(411774, 32);
    constant CORDIC_INV_GAIN : signed(31 downto 0) := to_signed(56281, 32);
    constant CORDIC_LATENCY : integer := 40;

    signal cordic_x_slv, cordic_y_slv : std_logic_vector(31 downto 0);
    signal cordic_mag_slv, cordic_pha_slv : std_logic_vector(31 downto 0);
    signal envelope_raw, envelope_s, phase_raw : data_t;
    signal phase_high_res : signed(31 downto 0);
    signal env_mult_long : signed(63 downto 0);
    signal valid_sr : std_logic_vector(CORDIC_LATENCY downto 0)
        := (others => '0');
    signal buffers_enable_int : std_logic := '0';
    signal last_phase_wrapped, last_phase_unwrapped : data_t
        := (others => '0');
    signal theta_unwrapped_reg, theta_norm_reg : data_t
        := (others => '0');

begin

    cordic_x_slv <= std_logic_vector(i_in);
    cordic_y_slv <= std_logic_vector(q_in);

    CORDIC_INST : cordic_rect_to_polar
    port map (
        clk => clk, x_in => cordic_x_slv, y_in => cordic_y_slv,
        x_out => cordic_mag_slv, phase_out => cordic_pha_slv
    );

    -- Correcao de ganho do envelope
    envelope_raw <= signed(cordic_mag_slv);
    env_mult_long <= envelope_raw * CORDIC_INV_GAIN;
    envelope_s <= env_mult_long(47 downto 16);

    -- Escala de fase
    phase_high_res <= signed(cordic_pha_slv);
    phase_raw <= resize(shift_right(phase_high_res, 13), 32);

    -- Linha de atraso para valid
    process(clk)
    begin
        if rising_edge(clk) then
            if rst = '1' then valid_sr <= (others => '0');
            else valid_sr <= valid_sr(CORDIC_LATENCY-1 downto 0)
                & input_valid; end if;
        end if;
    end process;

    buffers_enable_int <= valid_sr(CORDIC_LATENCY);
    buffers_enable     <= buffers_enable_int;

    -- Logica de fase (unwrap)
    process(clk)
        variable diff : data_t;
    begin
        if rising_edge(clk) then
            if rst = '1' then
                last_phase_wrapped <= (others => '0');
                last_phase_unwrapped <= (others => '0');
                theta_unwrapped_reg <= (others => '0');
                theta_norm_reg <= (others => '0');
                envelope_out <= (others => '0');
                phase_wrapped_out <= (others => '0');
            elsif buffers_enable_int = '1' then
                diff := phase_raw - last_phase_wrapped;
                if diff > PI_FIXED then
                    diff := diff - TWO_PI;
                elsif diff < -PI_FIXED then
                    diff := diff + TWO_PI;
                end if;

                theta_unwrapped_reg <= last_phase_unwrapped + diff;
                last_phase_wrapped <= phase_raw;
                last_phase_unwrapped <= last_phase_unwrapped + diff;
                theta_norm_out <= diff;
                envelope_out   <= envelope_s;
                phase_wrapped_out <= phase_raw;
            end if;
        end if;
    end process;

    theta_unwrapped_out <= theta_unwrapped_reg;

end Behavioral;
\end{lstlisting}


\subsection{Explica\c{c}\~{a}o do C\'{o}digo}

\subsubsection{Constantes}

\begin{itemize}
    \item \texttt{PI\_FIXED} = 205887 $\approx \pi \times 2^{16}$: Valor de $\pi$ em formato Q16.16.
    \item \texttt{TWO\_PI} = 411774 $\approx 2\pi \times 2^{16}$: Usado na l\'{o}gica de \textit{wrap}.
    \item \texttt{CORDIC\_INV\_GAIN} = 56281 $\approx K^{-1} \times 2^{16} = 0{,}8588 \times 65536$: Fator de corre\c{c}\~{a}o de ganho.
    \item \texttt{CORDIC\_LATENCY} = 40: N\'{u}mero de est\'{a}gios do pipeline CORDIC.
\end{itemize}

\subsubsection{Convers\~{a}o CORDIC}

O IP Core \texttt{cordic\_rect\_to\_polar} recebe as componentes $i_{in}$ e $q_{in}$ e produz:
\begin{itemize}
    \item \texttt{x\_out}: magnitude com ganho $K$ do CORDIC.
    \item \texttt{phase\_out}: fase em alta resolu\c{c}\~{a}o.
\end{itemize}

\subsubsection{Corre\c{c}\~{a}o de Ganho}

A magnitude \'{e} corrigida multiplicando por $K^{-1}$:
\begin{equation}
    \texttt{envelope\_s} = \left(\texttt{envelope\_raw} \times K^{-1}\right)[47\!:\!16]
\end{equation}
O produto de 64 bits \'{e} truncado para 32 bits extraindo os bits $[47\!:\!16]$, o que equivale a uma divis\~{a}o por $2^{16}$ (remo\c{c}\~{a}o da parte fracion\'{a}ria extra).

\subsubsection{Escala de Fase}

A sa\'{i}da de fase do CORDIC est\'{a} em um formato de alta resolu\c{c}\~{a}o e precisa ser escalonada:
\begin{equation}
    \texttt{phase\_raw} = \texttt{phase\_high\_res} \gg 13
\end{equation}

\subsubsection{Linha de Atraso para Valid}

Um registrador de deslocamento de 41 bits propaga o sinal \texttt{input\_valid} com exatamente 40 ciclos de atraso, sincronizando-o com a sa\'{i}da do CORDIC.

\subsubsection{Desembrulhamento de Fase}

A l\'{o}gica de \textit{phase unwrapping} garante continuidade da fase:
\begin{align}
    \Delta\theta(n) &= \theta_{raw}(n) - \theta_{raw}(n-1), \text{ ajustado para } [-\pi, +\pi] \label{eq:diff}\\
    \theta_{unwrapped}(n) &= \theta_{unwrapped}(n-1) + \Delta\theta(n) \label{eq:unwrap}
\end{align}

Se $\Delta\theta > \pi$, subtrai-se $2\pi$; se $\Delta\theta < -\pi$, soma-se $2\pi$. Isso evita descontinuidades quando a fase cruza $\pm\pi$.


% =======================================================================
\section{Testbench}
\label{sec:testbench}
% =======================================================================

O testbench \texttt{Input\_Processing\_Block\_tb} valida o bloco atrav\'{e}s de tr\^{e}s casos de teste:

\begin{enumerate}
    \item \textbf{Teste 1 -- Ganho CORDIC e Lat\^{e}ncia}: Envia $i_{in} = 1{,}0$, $q_{in} = 0{,}0$. Verifica se o envelope \'{e} $\approx 65536$ (ou seja, $1{,}0$ em Q16.16) e a fase $\approx 0$.
    \item \textbf{Teste 2 -- Fase $\pi$}: Envia $i_{in} = -1{,}0$, $q_{in} = 0{,}0$. Verifica se $\theta \approx 205887$ ($\pi$ em Q16.16).
    \item \textbf{Teste 3 -- Desembrulhamento de Fase}: Envia $i_{in} = -1{,}0$, $q_{in} = -0{,}1$. A fase passa ligeiramente al\'{e}m de $-\pi$, e a l\'{o}gica de \textit{unwrap} deve manter continuidade, resultando em $\theta_{unwrapped} > 210000$.
\end{enumerate}

\subsection{C\'{o}digo do Testbench}

\begin{lstlisting}[style=vhdlstyle, caption={Testbench do \texttt{Input\_Processing\_Block} (\texttt{input\_processing\_tb.vhd}).}, label={lst:tb}]
library IEEE;
use IEEE.STD_LOGIC_1164.ALL;
use IEEE.NUMERIC_STD.ALL;
use work.pkg_neural_types.ALL;

entity Input_Processing_Block_tb is
end Input_Processing_Block_tb;

architecture Behavioral of Input_Processing_Block_tb is

    component Input_Processing_Block
    Port (
        clk                : in  std_logic;
        rst                : in  std_logic;
        i_in               : in  data_t;
        q_in               : in  data_t;
        input_valid        : in  std_logic;
        envelope_out       : out data_t;
        theta_unwrapped_out: out data_t;
        theta_norm_out     : out data_t;
        buffers_enable     : out std_logic
    );
    end component;

    signal clk : std_logic := '0';
    signal rst : std_logic := '0';
    signal i_in : data_t := (others => '0');
    signal q_in : data_t := (others => '0');
    signal input_valid : std_logic := '0';

    signal envelope_out : data_t;
    signal theta_unwrapped_out : data_t;
    signal theta_norm_out : data_t;
    signal buffers_enable : std_logic;

    constant clk_period : time := 10 ns;

    -- Constantes Q16.16
    constant VAL_ONE : data_t := to_signed(65536, 32);
    constant VAL_NEG_ONE : data_t := to_signed(-65536, 32);
    constant VAL_HALF : data_t := to_signed(32768, 32);
    constant PI_VAL : integer := 205887;

begin

    uut: Input_Processing_Block PORT MAP (
        clk => clk, rst => rst,
        i_in => i_in, q_in => q_in,
        input_valid => input_valid,
        envelope_out => envelope_out,
        theta_unwrapped_out => theta_unwrapped_out,
        theta_norm_out => theta_norm_out,
        buffers_enable => buffers_enable
    );

    clk_process: process
    begin
        clk <= '0'; wait for clk_period/2;
        clk <= '1'; wait for clk_period/2;
    end process;

    stim_proc: process
        variable env_int : integer;
        variable theta_int : integer;
    begin
        -- Reset
        rst <= '1';
        wait for 100 ns;
        rst <= '0';
        wait for clk_period*5;

        report "--- TESTE 1: Ganho CORDIC e Latencia ---";
        i_in <= VAL_ONE;
        q_in <= (others => '0');
        input_valid <= '1';
        wait for clk_period;
        input_valid <= '0';

        wait until buffers_enable = '1';
        wait for clk_period * 2;

        env_int := to_integer(envelope_out);
        theta_int := to_integer(theta_unwrapped_out);

        report "Case 1 Output -> Env: "
            & integer'image(env_int)
            & " Theta: " & integer'image(theta_int);

        if abs(env_int - 65536) < 100 then
            report "PASS: Envelope Gain Correct";
        else
            report "FAIL: Envelope Gain Incorrect"
                severity error;
        end if;

        assert abs(theta_int) < 100
            report "FAIL: Phase 0 Incorrect"
            severity error;

        -- TESTE 2: Fase PI
        report "--- TESTE 2: Fase PI ---";
        wait for clk_period * 10;

        i_in <= VAL_NEG_ONE;
        q_in <= (others => '0');
        input_valid <= '1';
        wait for clk_period;
        input_valid <= '0';

        wait until buffers_enable = '1';
        wait for clk_period * 2;

        theta_int := to_integer(theta_unwrapped_out);
        report "Case 2 Output -> Theta: "
            & integer'image(theta_int);

        if abs(theta_int - PI_VAL) < 1000 then
            report "PASS: Phase PI Correct";
        else
            report "FAIL: Phase PI Incorrect"
                severity error;
        end if;

        -- TESTE 3: Fase Wrap
        report "--- TESTE 3: Fase Wrap ---";
        wait for clk_period * 10;

        i_in <= VAL_NEG_ONE;
        q_in <= to_signed(-6553, 32);
        input_valid <= '1';
        wait for clk_period;
        input_valid <= '0';

        wait until buffers_enable = '1';
        wait for clk_period * 2;

        theta_int := to_integer(theta_unwrapped_out);

        report "Case 3 Output -> Theta Unwrapped: "
            & integer'image(theta_int);

        if theta_int > 210000 then
             report "PASS: Phase Unwrap Logic";
        else
             report "FAIL: Phase Wrap failed"
                 severity error;
        end if;

        report "--- FIM DO TESTE ---";
        wait;
    end process;

end Behavioral;
\end{lstlisting}


% =======================================================================
\section{Resultados de Simula\c{c}\~{a}o}
\label{sec:resultados}
% =======================================================================

A simula\c{c}\~{a}o foi realizada utilizando o ISim (Xilinx ISE Design Suite). Todos os tr\^{e}s testes passaram com sucesso.

\subsection{Formas de Onda}

% NOTA: Substituir pelo caminho real da imagem de simulacao do ISim
% Para capturar: no ISim, File -> Export -> Export as Image (PNG)
\begin{figure}[H]
\centering
\fbox{\parbox{0.9\textwidth}{\centering\vspace{3cm}
    {\sffamily\large\color{black!50} [Inserir captura de tela do ISim aqui]}\\[5pt]
    {\sffamily\small\color{black!50} File $\rightarrow$ Export $\rightarrow$ Export as Image (PNG)}\\[5pt]
    {\sffamily\small\color{black!50} Salvar como: \texttt{isim\_input\_processing.png}}
\vspace{3cm}}}
\caption{Formas de onda da simula\c{c}\~{a}o do \texttt{Input\_Processing\_Block} no ISim.}
\label{fig:waveform}
\end{figure}

% Quando a imagem estiver disponivel, descomentar as linhas abaixo:
% \begin{figure}[H]
% \centering
% \includegraphics[width=\textwidth]{isim_input_processing.png}
% \caption{Formas de onda da simulacao do \texttt{Input\_Processing\_Block} no ISim.}
% \label{fig:waveform}
% \end{figure}


\subsection{Log do Simulador}

O log completo do ISim confirma que todos os testes foram aprovados:

\begin{lstlisting}[style=logstyle, caption={Log de sa\'{i}da do ISim para o testbench do \texttt{Input\_Processing\_Block}.}, label={lst:log}]
Finished circuit initialization process.
at 150 ns: Note: --- TESTE 1: Ganho CORDIC e Latencia ---
    (/input_processing_block_tb/).
at 575 ns: Note: Case 1 Output -> Env: 65536 Theta: 0
    (/input_processing_block_tb/).
at 575 ns: Note: PASS: Envelope Gain Correct
    (/input_processing_block_tb/).
at 575 ns: Note: --- TESTE 2: Fase PI ---
    (/input_processing_block_tb/).

# run 10.00us
at 1095 ns: Note: Case 2 Output -> Theta: 205886
    (/input_processing_block_tb/).
at 1095 ns: Note: PASS: Phase PI Correct
    (/input_processing_block_tb/).
at 1095 ns: Note: --- TESTE 3: Fase Wrap ---
    (/input_processing_block_tb/).
at 1615 ns: Note: Case 3 Output -> Theta Unwrapped: 212417
    (/input_processing_block_tb/).
at 1615 ns: Note: PASS: Phase Unwrap Logic
    (/input_processing_block_tb/).
at 1615 ns: Note: --- FIM DO TESTE ---
    (/input_processing_block_tb/)
\end{lstlisting}


\subsection{An\'{a}lise dos Resultados}

\begin{table}[H]
\centering
\caption{Resultados quantitativos dos testes de simula\c{c}\~{a}o.}
\label{tab:resultados}
\small
\begin{tabular}{llccc}
\hline
\textbf{Teste} & \textbf{Grandeza} & \textbf{Esperado} & \textbf{Obtido} & \textbf{Erro} \\
\hline
1 -- Ganho CORDIC & Envelope ($A$) & 65536 (1,0) & 65536 & 0 \\
1 -- Lat\^{e}ncia   & $\theta$ & 0 & 0 & 0 \\
2 -- Fase $\pi$      & $\theta_{unwrapped}$ & 205887 ($\pi$) & 205886 & $-1$ LSB \\
3 -- Phase Unwrap    & $\theta_{unwrapped}$ & $>$210000 & 212417 & -- \\
\hline
\end{tabular}
\end{table}

\begin{itemize}
    \item \textbf{Teste 1}: O envelope retornou exatamente $65536 = 1{,}0$ em Q16.16, confirmando que a corre\c{c}\~{a}o de ganho $K^{-1}$ funciona corretamente. A lat\^{e}ncia do CORDIC (40 ciclos) tamb\'{e}m foi verificada pelo atraso entre \texttt{input\_valid} e \texttt{buffers\_enable}.

    \item \textbf{Teste 2}: A fase para o vetor $(-1, 0)$ retornou $205886$, com erro de apenas 1~LSB em rela\c{c}\~{a}o ao valor te\'{o}rico $\pi \approx 205887$. Esse erro \'{e} aceit\'{a}vel para a precis\~{a}o Q16.16.

    \item \textbf{Teste 3}: Ap\'{o}s enviar o vetor $(-1, -0{,}1)$, cuja fase real \'{e} ligeiramente al\'{e}m de $-\pi$ (no terceiro quadrante), a l\'{o}gica de \textit{unwrap} detectou corretamente o salto de $2\pi$ e acumulou o valor $\theta_{unwrapped} = 212417$, mantendo continuidade.
\end{itemize}


% =======================================================================
\section{Conclus\~{a}o e Pr\'{o}ximos Passos}
% =======================================================================

O \texttt{Input\_Processing\_Block} foi implementado e validado com sucesso. Os tr\^{e}s testes confirmaram:

\begin{itemize}
    \item Corre\c{c}\~{a}o do ganho CORDIC com precis\~{a}o exata.
    \item C\'{a}lculo correto de fase com erro $\leq 1$~LSB.
    \item Funcionamento da l\'{o}gica de desembrulhamento de fase.
    \item Lat\^{e}ncia correta de 40 ciclos propagada pelo \texttt{valid\_sr}.
\end{itemize}

\noindent\textbf{Pr\'{o}ximos passos:}
\begin{enumerate}
    \item Implementar e validar as LUTs trigonom\'{e}tricas (\texttt{LUT\_Trig}) com pipeline de 4 est\'{a}gios.
    \item Implementar os registradores de deslocamento (\texttt{Shift\_Register}).
    \item Integrar os blocos no wrapper do Est\'{a}gio~1 (\texttt{Top\_NeuralNet\_S1}).
    \item Validar o Est\'{a}gio~1 completo antes de prosseguir para a camada convolucional.
\end{enumerate}

\end{document}
